\subsection{Simultaneous predecessor-bit folds}
\label{sec:simultaneous-folds}

The following finite certificate criterion is sufficient for the reader-state
folds used in our search.  It is stronger than checking each proposed fold in
isolation: a helper used by one fold may itself be changed by another fold.

Let $M$ be the original binary transition table.  A \emph{pure reader} $q$ has
rows $M(q,b)=(b,R,q_b)$ for $b\in\{0,1\}$.  For each state $q$, conservatively
collect the possible predecessor bits from every incoming transition: an
incoming right-moving row that writes $a$ contributes $a$, and every incoming
left-moving row contributes both bits.  Suppose $q$ has predecessor set
$\{0\}$ and $r$ has predecessor set $\{1\}$.  The two states can potentially
share one state, because the tape cell immediately left of the head records
which of the two original states was entered.

A proposed fold retains $q$, aliases $r$ to $q$, and, for each scanned bit $b$,
either retains their common row or writes $b$ and moves left to an existing
helper $h_b$.  The usual three-step helper satisfies
\[
 M(h_b,0)=(0,R,a),\qquad M(h_b,1)=(1,R,c),\qquad
 M(a,b)=M(q,b),\qquad M(c,b)=M(r,b).
\]
These equations suffice only if all rows in this expansion retain their stated
behavior.  Choosing a matching of pairs without checking this condition can
produce invalid simultaneous folds.

Our checker therefore applies \emph{all} selected folds first.  Write
$\mu(r)=q$ for each removed state and let $M'$ be the resulting complete table,
with every transition target renamed by $\mu$.  For each selected pair
$(q,r)$, each $j,b\in\{0,1\}$, let $s_0=q$ and $s_1=r$.  The checker executes
$M'$ from state $\mu(s_j)$, head position $0$, and the two known tape cells
\[
 T(-1)=j,\qquad T(0)=b.
\]
It requires a finite positive execution ending with the head at $1$, state
$\mu((M(s_j,b))_{\rm next})$, and both known cells unchanged.  Reading a cell
outside $\{-1,0\}$ before that endpoint, encountering an undefined row, or
repeating a configuration rejects the proposal.  The implementation uses a
finite fuel bound; exhausting it also rejects, so it can lose opportunities
but cannot certify an unverified expansion.  Recorded successful traces give
finite certificates that can be checked independently.

\paragraph{Soundness.}
At an original entry to $q$ or $r$, the incoming transition has written the
required predecessor bit, since any left-moving incoming transition would
have excluded that state from the candidate sets.  Thus the local certificate
applies independently of the tape outside the two inspected cells.  Unfolded
states take the same one step, with only the destination renamed; folded
states take their certified positive finite sequence.  Induction on the
original run therefore produces a simulation in $M'$, preserving tape and
head at every original step boundary.  Every inserted sequence terminates
and has no intermediate halting configuration.  Conversely, if $M'$ halted
before an original halting boundary, finitely many of these positive
sequences would extend past that time, contradicting determinism.  Hence the
machines have the same halting behavior from any common start satisfying the
incoming-bit invariant.  The start state is excluded from the folds.

All reasoning here is about fixed finite tables and fixed finite local traces;
the run induction and the bounded reflection argument are formalizable in
$\mathsf{PA}$.  This is a mathematical certificate argument, not a claim that
our Python implementation has itself been verified by the Lean kernel.
